% v11 staging module. General two-high sign theorem and triangular-support obstruction.
% Assumes kl10:two-high-parametrization, kl10:two-high-positive,
% and the already established opposing-end cancellation of D.

\section{二高分母の符号と三角形成分の接着障害}
\label{ls11:section}

二高分母の factorization 表示の正値性条件は，分母が異なる場合にも
単純な符号条件まで強められる．これを共有 puncture の整数 clutching と組み合わせると，
高分母 support の一つの連結成分が三角形になる場合を，分母の上限なしで排除できる．

\begin{theorem}[二高分母の正値性と符号]\label{ls11:two-high-signs}
命題~\ref{kl10:two-high-parametrization} の標準表示に従う
分母 $(a,b,1,1)$，$a,b>1$，の整数 numerator matrix を考える．
正の局所 weights が存在することは
\[
 u<0<v
\]
と同値である．従って，正の局所模型の各高分母 row は，もう一つの高分母へ負の entry を持ち，
二つの低分母へは正と負の entry を一つずつ持つ．
この結論は $a=b$ を仮定しない．
\end{theorem}

\begin{proof}
標準表示の記号を用いる．$A=a+x+u\geq x+1\geq2$ であり，
pair-product 条件と $x=(a-1)/r$，$y=(b-1)/(r+1)$ から
\begin{equation}\label{ls11:linear-v}
 A(r+1)v=(1-ru)b-(A+r+1)
\end{equation}
を得る．$u<0$ なら $1-ru>0$ である．もし $v\leq-1$ なら
\[
 (1-ru)b=A(r+1)v+A+r+1
 \leq r+1-Ar\leq0
\]
となり，$b>1$ に反する．$v\ne0$ なので $v>0$ である．
また，$u<0$ なら既知の正値性条件~\eqref{kl10:two-high-positive} が成立する．

残る $u>0$ の場合が正の weights を持たないことを示す．
この場合は $v<0$ なので $v=-z$，$z>0$ と書く．
\[
 S=\frac{a-1}{u},\qquad T=\frac{b-1}{z},\qquad
 d=T-S-1
\]
と置くと $S,T$ は正の整数である．正値性は
$rud<a=uS+1$ を要求し，整数性から
\begin{equation}\label{ls11:cycle-bound}
 rd\leq S
\end{equation}
となる．次の正の量と整数を導入する：
\[
 \begin{split}
 F&=(2r+1)S+r-r^2d,\qquad G=(r+1)(S+r),\\
 H&=r(r+1),\qquad K=r(r+T+1).
 \end{split}
\]
pair-product 条件を $a=uS+1$，$b=zT+1$，$T=S+d+1$ によって書き直すと
\begin{equation}\label{ls11:bilinear-positive}
 u(zF-G)=H-zK
\end{equation}
を得る．式~\eqref{ls11:cycle-bound} より
\[
 F\geq(r+1)S+r>0.
\]
一方，$z\geq1$，$T\geq1$ より
$H-zK\leq H-K=-rT<0$ である．従って $zF<G$ であり，特に
\[
 0<G-F=r\{r(d+1)-S\}.
\]
従って
\begin{equation}\label{ls11:short-interval}
 rd\leq S<r(d+1).
\end{equation}
$S>0$ なので $d\geq0$ である．

$d=0$ なら，二つの低分母頂点間の entries は両方とも零である．
これらの頂点を交換すると，再び同じ標準表示になり，新しい $r,S$ はそれぞれ $S,r$ となる．
式~\eqref{ls11:short-interval} は元の標識で $S<r$，交換後には $r<S$ を与えるので矛盾する．

$d\geq1$ なら $S\geq r$ であり，
\[
 1\leq z<\frac GF
 \leq\frac{(r+1)(S+r)}{(r+1)S+r}
 =1+\frac{r^2}{(r+1)S+r}<2.
\]
従って $z=1$ である．$S=rd+e$，$0\leq e<r$ と書けば
\[
 T=S+d+1=(r+1)d+e+1.
\]
標準表示の整除条件 $r+1\mid b-1=zT$ は，$z=1$ より $r+1\mid e+1$ を要求する．
しかし $1\leq e+1\leq r$ なので矛盾する．これで $u>0$ は排除された．
\end{proof}

この定理は，既知の factorization 表示における cycle 不等式を削除するものではなく，
その整数条件を全て用いて，正値な場合を $u<0<v$ にちょうど絞り込む結果である．
局所の正の chart の実現は命題~\ref{q9:four-line-criterion} に従う．

\begin{lemma}[両端が低分母に接する puncture]\label{ls11:low-puncture}
共有 branch の分母を $q>1$ とし，その同じ marked puncture が，
両側の core において分母一の ray に向かうとする．
共有 branch からその ray へ向かう二つの numerators を $p,p'$ と書く．
整数 clutching $q\mid p+p'$ が成立するなら，次のいずれかに限られる：
\[
 p'=-p,\qquad
 \{p,p'\}=\{1,q-1\},\qquad
 \{p,p'\}=\{-1,-(q-1)\}.
\]
この三場合のいずれでも，逆向きの numerators の和は元の和と等しい：
\[
 \frac{q-1}{p}+\frac{q-1}{p'}=p+p'.
\]
meridian node 表示では，その puncture の node 数はそれぞれ $1,0,2$ である．
\end{lemma}

\begin{proof}
crossing determinant により $p,p'$ は $q-1$ の非零 signed divisors である．
従って $|p+p'|<2q$ であり，その和は $-q,0,q$ のいずれかとなる．
和が零なら第一の場合である．和が $q$ なら両者は正であり，
両者とも $q-1$ より小さい divisor なら各々 $(q-1)/2$ 以下なので矛盾する．
従って一方が $q-1$，他方が一である．負の場合も同様である．
逆向きの和についての等式と node 数は，各場合への直接代入で従う．
\end{proof}

\begin{theorem}[高分母の三角形成分の不存在]\label{ls11:no-triangle-component}
現在の全平面五次の構成条件の下で，大域作用適合 atlas が存在するとする．
その高分母 support
\[
 H=\bigl\{\{i,j\}:q_{ij}>1\bigr\}\subset K_5
\]
は，三角形を一つの連結成分として持たない．
これは分母の大きさによらない一般障害である．
証明に用いるのは局所の正値性，row sums，crossing determinants，
marked puncture の整数 clutching，および対向端の $D$ の相殺である．
node 数の非負性や各直線の全 node 数五は，この定理には必要ない．
\end{theorem}

\begin{proof}
三角形の core labels を $0,1,2$ とし，添字は三を法として読む．
辺 $\{i,i+1\}$ の分母を $q_i>1$ と書く．
三角形成分という仮定から，この三つの cores はそれぞれ二高分母を持ち，
残る二つの core labels へ向かう分母は一である．
各共有 branch の二つの外側 punctures は，それら二つの固定された labels に対応する．

定理~\ref{ls11:two-high-signs} より，その二つの outgoing entries は各端で一正一負である．
同じ marked punctures を接続するとき，符号が対応するなら補題~\ref{ls11:low-puncture} により
二つの和は $q_i,-q_i$，符号が逆に対応するなら零，零となる．
従って二つの外側 punctures の outgoing sums の総和は零であり，
同補題により incoming sums の総和も零である．

core $i$ における二つの高分母 rays の numerators を
\[
 p_{i;i-1,i+1}=-\alpha_i,\qquad
 p_{i;i+1,i-1}=-\beta_i,
 \qquad \alpha_i,\beta_i>0
\]
と書く．局所 crossing determinant は
\begin{equation}\label{ls11:triangle-products}
 \alpha_i\beta_i=q_{i-1}q_i-1
\end{equation}
である．両端の outgoing row sums から，辺 $\{i,i+1\}$ について
\[
 \beta_i+\alpha_{i+1}=2q_i
\]
となる．また対向端の $D$ の相殺と incoming sums の総和零から
\[
 \alpha_i+\beta_{i+1}=2q_i
\]
を得る．差を取ると $\alpha_i-\beta_i$ は $i$ によらない整数 $\delta$ となる．
meridian node 表示を用いる場合には，最初の等式は三角形に向かう puncture の node 数三を意味するが，
ここではこの個数や全 node budget を別の仮定として用いていない．

$T_i=(\alpha_i+\beta_i)/2>0$ と置くと，
\[
 q_i=\frac{T_i+T_{i+1}}2,
 \qquad
 (T_{i-1}+T_i)(T_i+T_{i+1})=4T_i^2-\delta^2+4.
\]
$T_i$ が最大となる頂点では左辺は $4T_i^2$ 以下なので $\delta^2\geq4$，
最小となる頂点では逆に $\delta^2\leq4$ である．
従って $\delta^2=4$ であり，最大頂点での等号から三つの $T_i$ は全て等しい．
従って三辺の分母も全て同じ $q$ となり，各 core の二つの high--high numerators は
$-q-1,-q+1$ となる．

この特別な二高分母模型を標準表示で調べる．
$m=q-1=t r(r+1)$，$x=t(r+1)$，$y=tr$ とし，$u=-s<0$，$v=z>0$ と置く．
high--high entries は $-A,-B$，
\[
 A=q+x-s,\qquad B=q-y+z,\qquad \{A,B\}=\{q-1,q+1\}
\]
である．もし $A=q-1$ なら $s=x+1$ であり，$s\mid m=rx$ から $x+1\mid r$ となる．
しかし $x=t(r+1)\geq r+1$ なので不可能である．
従って $A=q+1$ であり，$s=x-1$，$z=y-1$ となる．
$s\mid rx$ より $x-1\mid r$，一方 $x-1\geq r$ なので $t=1$，$s=r$ である．
さらに $z=r-1>0$，$z\mid r(r+1)$ より $r-1\mid2$ となる．
従って $r=2,3$，すなわち $q=7,13$ に限られる．

直接計算すると，その二高分母における $D$ の組は，順序を除き
\[
 q=7:\quad (0,-2),\qquad
 q=13:\quad (0,-1)
\]
である．どちらも高分母二頂点の係数和は負である．
三つの cores にわたりこの和を足すと負になるが，各三角形辺の対向端では相殺するため零でなければならない．
この矛盾が結論を与える．
\end{proof}

本定理は，高分母 support が三角形を部分 graph として含む場合を全て排除するものではない．
三つの cores のいずれかにさらに高分母 branch が接続する場合には，
二高分母の符号定理をその core へ適用できない．
また，四角形・五角形や一般の path 成分の排除は，本定理だけからは従わない．
